% ----------------------------------------------------------------------
% Introdução: fica antes da primeira parte
% ----------------------------------------------------------------------
\chapter{Introdução}
\label{cap:introducao}

\instrucao{De cinco a dez páginas. Na tese, a introdução apresenta o tema, o problema, a tese defendida (a afirmação que o trabalho sustenta), os objetivos, a justificativa e o percurso. Termine com um mapa das partes.}

Os repositórios institucionais tornaram-se o principal canal de acesso aberto à produção das universidades brasileiras. Para que seus registros sejam coletados por agregadores, como a Biblioteca Digital Brasileira de Teses e Dissertações (BDTD), os metadados precisam seguir padrões comuns e ser expostos por protocolos como o \textit{Open Archives Initiative Protocol for Metadata Harvesting} (OAI-PMH) \cite{lagoze2001}.

\section{Problema e tese}
\instrucao{Formule o problema como pergunta e, logo depois, a tese como resposta provisória que o trabalho vai demonstrar.}

O problema desta pesquisa é: em que medida os repositórios institucionais brasileiros são interoperáveis quanto aos seus metadados? Defende-se a tese de que a interoperabilidade sintática está consolidada, mas a semântica permanece frágil, porque o controle de autoridades não acompanhou a expansão dos repositórios.

\section{Objetivos}
\instrucao{Um objetivo geral e de três a cinco específicos, cada um ligado a um capítulo.}
O objetivo geral é propor e aplicar um modelo de avaliação da interoperabilidade de metadados em repositórios institucionais. Os objetivos específicos são:
\begin{enumerate}
  \item sistematizar as dimensões de interoperabilidade descritas na literatura;
  \item construir indicadores mensuráveis para cada dimensão;
  \item aplicar o modelo a uma amostra de repositórios brasileiros;
  \item propor uma agenda priorizada de melhorias.
\end{enumerate}

\section{Organização da tese}
A Parte~\ref{part:fundamentos} reúne os fundamentos teóricos, nos Capítulos~\ref{cap:metadados} e~\ref{cap:interoperabilidade}. A Parte~\ref{part:empirica} descreve o método e os resultados, nos Capítulos~\ref{cap:metodo} e~\ref{cap:resultados}. O Capítulo~\ref{cap:conclusao} retoma a tese à luz dos achados.
